\section{Approach}
\label{sec:approach}

\subsection{Delegation edges}

An edge $p \rightarrow q$ means $p$ can cause code to run with $q$'s
permissions, so the edges are what produce the $E(p)$ of
Equation~\ref{eq:eff}. Three rules produce them over the deployment
$D = (P, A, G, T)$ of Section~\ref{sec:model}, and each is a conjunction,
which matters for Section~\ref{sec:blocked}.
%
\begin{align}
\label{eq:assume}
p \xrightarrow{\text{assume}} q \;&\iff\;
  (p, \iam{sts:AssumeRole}, q) \in G \,\wedge\, (p,q) \in T \\
\label{eq:pass}
p \xrightarrow{\text{pass}} q \;&\iff\;
  (p, \iam{iam:PassRole}, q) \in G \,\wedge\, \nonumber\\
  &\qquad \exists\, \sigma \in \Sigma :\, (p, \sigma, \cdot) \in G \\
\label{eq:replace}
p \xrightarrow{\text{replace}} \rho(a) &\iff
  (p, \api{UpdateAgentRuntime}, a) \in G \wedge \nonumber\\
  &\qquad (p, \iam{iam:PassRole}, \rho(a)) \in G
\end{align}
%
$\Sigma$ is the set of actions that start a workload under a passed role.
Equation~\ref{eq:replace} is the agent specific rule. The platform's update
API requires the role in the same call as the artifact, so a caller holding the
first conjunct and not the second cannot make the call.

Figure~\ref{fig:mechanism} is Equation~\ref{eq:replace} drawn: the gate on the
middle arrow is the second conjunct, and the dashed box beside it is the
condition under which the managed policy satisfies it.

\subsection{Computing the closure}

Algorithm~\ref{alg:reach} computes the fixpoint of Equation~\ref{eq:eff} and
the difference of Equation~\ref{eq:esc}. The outer loop visits each principal
once and the inner worklist follows edges. A visited set keeps a cycle from
diverging, such as two roles that may assume each other. Each principal's
closure is memoised, so the edge set is walked once per principal and the cost
is $O(|P| \cdot (|P| + |G|))$.

\begin{algorithm}[t]
\caption{Escalation by closure over delegation edges}
\label{alg:reach}
\begin{algorithmic}[1]
\Function{Eff}{$p$, $D$}
  \If{$p \in \mathit{memo}$} \Return $\mathit{memo}[p]$ \EndIf
  \State $\mathit{memo}[p] \gets \auth(p)$ \Comment{guards cycles}
  \State $S \gets \auth(p)$
  \ForAll{$q \in E(p)$} \Comment{Eq.~\ref{eq:assume}--\ref{eq:replace}}
    \State $S \gets S \cup \Call{Eff}{q, D}$
  \EndFor
  \State $\mathit{memo}[p] \gets S$
  \State \Return $S$
\EndFunction
\Statex
\Function{Findings}{$D$}
  \State $F \gets \emptyset$
  \ForAll{$p \in P$}
    \State $\Delta \gets \Call{Eff}{p, D} \setminus \auth(p)$
    \If{$\Delta \neq \emptyset$}
      \State $F \gets F \cup \{(p, \Delta, \Call{Witness}{p, \Delta})\}$
    \EndIf
    \ForAll{$(q, \alpha) \in \Call{Blocked}{p, D}$}
      \State record $p \rightarrow q$ blocked on $\alpha$
    \EndFor
  \EndFor
  \State \Return $F$
\EndFunction
\end{algorithmic}
\end{algorithm}

\textsc{Witness} returns the path that carried an element of $\Delta$ into
$\eff(p)$, so every finding can be checked against the source by hand.

\subsection{Complete and blocked}
\label{sec:blocked}

Each edge rule is a conjunction. \textsc{Blocked} evaluates the conjuncts
separately and reports the pairs where exactly one is false, naming the
missing action. A path is \emph{complete} when every conjunct of every edge on
it holds, and \emph{blocked} when one does not.

The two are reported separately. A blocked candidate names the single grant
that would make the deployment exploitable, which is a rule for the next change
review.

\subsection{What counts as an agent}

A resource is an agent when the deployment declares an agent runtime for the
platform. The analysis does not guess from a name or an image. A deployment
that runs several agents inside one runtime is therefore one agent here,
because that is what it declares and because the platform has no identity
boundary inside a runtime for the analysis to find.
